% --- Archon physics formalization source begin ---
% archon:physics
% archon:covers PhyXMiniProblems/problem_phyx_mini_0724.lean
% archon:source-report reports/phyx_mini/problem_phyx_mini_0724.source.json

\chapter{Physics problem phyx\_mini\_0724}
\label{ch:PhyXMiniProblems_problem_phyx_mini_0724}

\paragraph{Problem source.}
\#\# Physical scenario

Water flows through the pipes shown in figure. The water's speed through the lower pipe is \textbackslash{}(5.0\textbackslash{},\textbackslash{}mathrm\{m/s\}\textbackslash{}) and a pressure gauge reads \textbackslash{}(75\textbackslash{},\textbackslash{}mathrm\{kPa\}\textbackslash{}).

\#\# Question

What is the reading of the pressure gauge on the upper pipe?

\#\# Answer choices

- A: A: 4.4kPa

- B: B: 4.8kPa

- C: C: 4.6kPa

- D: D: 5.0kPa

\#\# Figure caption (auxiliary; use the image as primary evidence)

The image depicts a fluid dynamics scenario involving flow through a pipe that changes in diameter and elevation. Here’s a detailed description:

- **Flow Direction**: Indicated by green arrows, the flow proceeds from left to right.

- **Section 1**: 
  - The pipe has a diameter of 6.0 cm.
  - A pressure gauge shows 75 kPa.
  - The flow velocity is labeled as 5.0 m/s.
  - This section is marked with the number "1".

- **Section 2**: 
  - The pipe reduces in diameter to 4.0 cm.
  - There is an unknown pressure indicated by a question mark on the gauge.
  - The flow velocity is labeled as \textbackslash{}( v\_2 \textbackslash{}).
  - This section is marked with the number "2" and is 2.0 m higher in elevation compared to section 1.

The pipe has an upward slope between sections 1 and 2, indicating a change in height and possibly pressure and velocity between the two sections.

\#\# Dataset metadata

Dynamics; Multi-Formula Reasoning, Physical Model Grounding Reasoning

\paragraph{Recorded answer/context.}
C: C: 4.6kPa

\paragraph{Figure/image path.}
/root/proposal\_for\_physic/science-mango/phyx\_mini\_run/phyx\_data/test\_image/724.png

\paragraph{Formalization target.}
create a compiling Lean file with sorry bodies at `PhyXMiniProblems/problem\_phyx\_mini\_0724.lean`. The Lean declarations must preserve the physical quantities, dimensions or dimensional roles, figure labels, governing-law hypotheses, and final relation expressed by this problem.
Use Mathlib/Physlib names found through LeanExplore where available. If a physics API is missing, introduce faithful local abstractions rather than scalar placeholder aliases.

\begin{theorem}[Physics formalization target]
\label{thm:physics:phyx_mini_0724:target}
\lean{PhyXMiniProblems.ProblemPhyXMini0724.problem_phyx_mini_0724}
\uses{def:physics:phyx-mini-0724:phyxminiproblems-problemphyxmini0724-pressureinpascals, def:physics:phyx-mini-0724:phyxminiproblems-problemphyxmini0724-pressureinkilopascals, def:physics:phyx-mini-0724:phyxminiproblems-problemphyxmini0724-risingpipesetup, def:physics:phyx-mini-0724:phyxminiproblems-problemphyxmini0724-matchesidealwaterpipescenario, def:physics:phyx-mini-0724:phyxminiproblems-problemphyxmini0724-matchesprimarypipefigure, def:physics:phyx-mini-0724:phyxminiproblems-problemphyxmini0724-matchesproblemreadouts, def:physics:phyx-mini-0724:phyxminiproblems-problemphyxmini0724-usestextbookwaterdensityandgravity, def:physics:phyx-mini-0724:phyxminiproblems-problemphyxmini0724-satisfiescircularrisingpipegeometry, def:physics:phyx-mini-0724:phyxminiproblems-problemphyxmini0724-satisfiesidealpipeflowlaws, def:physics:phyx-mini-0724:phyxminiproblems-problemphyxmini0724-isuniquematchinguppergaugepressure}
The assigned autoformalize agent should translate this physics problem into Lean declarations in the covered file, with theorem and lemma proof bodies written as `by sorry`.
\end{theorem}
\begin{proof}
This is an autoformalization task, not a proof task. Produce statements that can later be proved by the physics prover without weakening the physical model.
\end{proof}

% --- Archon declaration topology begin ---
\section{Declaration topology}
\begin{definition}[mass Density Dimension]
  \label{def:physics:phyx-mini-0724:phyxminiproblems-problemphyxmini0724-massdensitydimension}
  \lean{PhyXMiniProblems.ProblemPhyXMini0724.massDensityDimension}
  The physical dimension of mass density, M L⁻³
\end{definition}
\begin{proof}
  This is the declared model definition of mass Density Dimension.
\end{proof}
\begin{definition}[acceleration Dimension]
  \label{def:physics:phyx-mini-0724:phyxminiproblems-problemphyxmini0724-accelerationdimension}
  \lean{PhyXMiniProblems.ProblemPhyXMini0724.accelerationDimension}
  The physical dimension of acceleration, L T⁻²
\end{definition}
\begin{proof}
  This is the declared model definition of acceleration Dimension.
\end{proof}
\begin{definition}[Length Quantity]
  \label{def:physics:phyx-mini-0724:phyxminiproblems-problemphyxmini0724-lengthquantity}
  \lean{PhyXMiniProblems.ProblemPhyXMini0724.LengthQuantity}
  A nonnegative physical length, independent of a choice of units
\end{definition}
\begin{proof}
  This is the declared model definition of Length Quantity.
\end{proof}
\begin{definition}[Mass Density Quantity]
  \label{def:physics:phyx-mini-0724:phyxminiproblems-problemphyxmini0724-massdensityquantity}
  \lean{PhyXMiniProblems.ProblemPhyXMini0724.MassDensityQuantity}
  \uses{def:physics:phyx-mini-0724:phyxminiproblems-problemphyxmini0724-massdensitydimension}
  A nonnegative physical mass density
\end{definition}
\begin{proof}
  This is the declared model definition of Mass Density Quantity.
\end{proof}
\begin{definition}[Acceleration Quantity]
  \label{def:physics:phyx-mini-0724:phyxminiproblems-problemphyxmini0724-accelerationquantity}
  \lean{PhyXMiniProblems.ProblemPhyXMini0724.AccelerationQuantity}
  \uses{def:physics:phyx-mini-0724:phyxminiproblems-problemphyxmini0724-accelerationdimension}
  A nonnegative physical acceleration magnitude
\end{definition}
\begin{proof}
  This is the declared model definition of Acceleration Quantity.
\end{proof}
\begin{definition}[length In Meters]
  \label{def:physics:phyx-mini-0724:phyxminiproblems-problemphyxmini0724-lengthinmeters}
  \lean{PhyXMiniProblems.ProblemPhyXMini0724.lengthInMeters}
  \uses{def:physics:phyx-mini-0724:phyxminiproblems-problemphyxmini0724-lengthquantity}
  Coherent-SI metre readout of a physical length
\end{definition}
\begin{proof}
  This is the declared model definition of length In Meters.
\end{proof}
\begin{definition}[length In Centimeters]
  \label{def:physics:phyx-mini-0724:phyxminiproblems-problemphyxmini0724-lengthincentimeters}
  \lean{PhyXMiniProblems.ProblemPhyXMini0724.lengthInCentimeters}
  \uses{def:physics:phyx-mini-0724:phyxminiproblems-problemphyxmini0724-lengthquantity, def:physics:phyx-mini-0724:phyxminiproblems-problemphyxmini0724-lengthinmeters}
  Centimetre readout used by the two pipe-diameter labels
\end{definition}
\begin{proof}
  This is the declared model definition of length In Centimeters.
\end{proof}
\begin{definition}[area In Square Meters]
  \label{def:physics:phyx-mini-0724:phyxminiproblems-problemphyxmini0724-areainsquaremeters}
  \lean{PhyXMiniProblems.ProblemPhyXMini0724.areaInSquareMeters}
  Coherent-SI square-metre readout of a physical area
\end{definition}
\begin{proof}
  This is the declared model definition of area In Square Meters.
\end{proof}
\begin{definition}[speed In Meters Per Second]
  \label{def:physics:phyx-mini-0724:phyxminiproblems-problemphyxmini0724-speedinmeterspersecond}
  \lean{PhyXMiniProblems.ProblemPhyXMini0724.speedInMetersPerSecond}
  Coherent-SI metre-per-second readout of a speed magnitude
\end{definition}
\begin{proof}
  This is the declared model definition of speed In Meters Per Second.
\end{proof}
\begin{definition}[pressure In Pascals]
  \label{def:physics:phyx-mini-0724:phyxminiproblems-problemphyxmini0724-pressureinpascals}
  \lean{PhyXMiniProblems.ProblemPhyXMini0724.pressureInPascals}
  Coherent-SI pascal readout of a pressure
\end{definition}
\begin{proof}
  This is the declared model definition of pressure In Pascals.
\end{proof}
\begin{definition}[pressure In Kilopascals]
  \label{def:physics:phyx-mini-0724:phyxminiproblems-problemphyxmini0724-pressureinkilopascals}
  \lean{PhyXMiniProblems.ProblemPhyXMini0724.pressureInKilopascals}
  \uses{def:physics:phyx-mini-0724:phyxminiproblems-problemphyxmini0724-pressureinpascals}
  Kilopascal readout used by both pressure gauges and all answer choices
\end{definition}
\begin{proof}
  This is the declared model definition of pressure In Kilopascals.
\end{proof}
\begin{definition}[density In Kilograms Per Cubic Meter]
  \label{def:physics:phyx-mini-0724:phyxminiproblems-problemphyxmini0724-densityinkilogramspercubicmeter}
  \lean{PhyXMiniProblems.ProblemPhyXMini0724.densityInKilogramsPerCubicMeter}
  \uses{def:physics:phyx-mini-0724:phyxminiproblems-problemphyxmini0724-massdensityquantity}
  Coherent-SI kilogram-per-cubic-metre readout of a mass density
\end{definition}
\begin{proof}
  This is the declared model definition of density In Kilograms Per Cubic Meter.
\end{proof}
\begin{definition}[acceleration In Meters Per Second Squared]
  \label{def:physics:phyx-mini-0724:phyxminiproblems-problemphyxmini0724-accelerationinmeterspersecondsquared}
  \lean{PhyXMiniProblems.ProblemPhyXMini0724.accelerationInMetersPerSecondSquared}
  \uses{def:physics:phyx-mini-0724:phyxminiproblems-problemphyxmini0724-accelerationquantity}
  Coherent-SI metre-per-second-squared readout of an acceleration
\end{definition}
\begin{proof}
  This is the declared model definition of acceleration In Meters Per Second Squared.
\end{proof}
\begin{definition}[Pipe Section]
  \label{def:physics:phyx-mini-0724:phyxminiproblems-problemphyxmini0724-pipesection}
  \lean{PhyXMiniProblems.ProblemPhyXMini0724.PipeSection}
  The numbered cross sections marked by black dots in the primary image
\end{definition}
\begin{proof}
  This is the declared model definition of Pipe Section.
\end{proof}
\begin{definition}[Fluid Kind]
  \label{def:physics:phyx-mini-0724:phyxminiproblems-problemphyxmini0724-fluidkind}
  \lean{PhyXMiniProblems.ProblemPhyXMini0724.FluidKind}
  Fluids relevant to the physical role of the pipe contents
\end{definition}
\begin{proof}
  This is the declared model definition of Fluid Kind.
\end{proof}
\begin{definition}[Pressure Reference]
  \label{def:physics:phyx-mini-0724:phyxminiproblems-problemphyxmini0724-pressurereference}
  \lean{PhyXMiniProblems.ProblemPhyXMini0724.PressureReference}
  The reference convention used by a pressure-gauge readout
\end{definition}
\begin{proof}
  This is the declared model definition of Pressure Reference.
\end{proof}
\begin{definition}[Horizontal Figure Location]
  \label{def:physics:phyx-mini-0724:phyxminiproblems-problemphyxmini0724-horizontalfigurelocation}
  \lean{PhyXMiniProblems.ProblemPhyXMini0724.HorizontalFigureLocation}
  Horizontal ordering of the two marked sections in the supplied raster
\end{definition}
\begin{proof}
  This is the declared model definition of Horizontal Figure Location.
\end{proof}
\begin{definition}[Vertical Figure Location]
  \label{def:physics:phyx-mini-0724:phyxminiproblems-problemphyxmini0724-verticalfigurelocation}
  \lean{PhyXMiniProblems.ProblemPhyXMini0724.VerticalFigureLocation}
  Vertical ordering of the two marked sections in the supplied raster
\end{definition}
\begin{proof}
  This is the declared model definition of Vertical Figure Location.
\end{proof}
\begin{definition}[Horizontal Flow Direction]
  \label{def:physics:phyx-mini-0724:phyxminiproblems-problemphyxmini0724-horizontalflowdirection}
  \lean{PhyXMiniProblems.ProblemPhyXMini0724.HorizontalFlowDirection}
  Direction of a green flow arrow in the supplied raster
\end{definition}
\begin{proof}
  This is the declared model definition of Horizontal Flow Direction.
\end{proof}
\begin{definition}[Water Pipe Figure]
  \label{def:physics:phyx-mini-0724:phyxminiproblems-problemphyxmini0724-waterpipefigure}
  \lean{PhyXMiniProblems.ProblemPhyXMini0724.WaterPipeFigure}
  \uses{def:physics:phyx-mini-0724:phyxminiproblems-problemphyxmini0724-pipesection, def:physics:phyx-mini-0724:phyxminiproblems-problemphyxmini0724-horizontalfigurelocation, def:physics:phyx-mini-0724:phyxminiproblems-problemphyxmini0724-verticalfigurelocation, def:physics:phyx-mini-0724:phyxminiproblems-problemphyxmini0724-horizontalflowdirection}
  Literal labels and qualitative features visible in image 724.png. The unknown upper gauge is represented by its question-mark label, not by a numerical pressure value
\end{definition}
\begin{proof}
  This is the declared model definition of Water Pipe Figure.
\end{proof}
\begin{definition}[Pipe Section State]
  \label{def:physics:phyx-mini-0724:phyxminiproblems-problemphyxmini0724-pipesectionstate}
  \lean{PhyXMiniProblems.ProblemPhyXMini0724.PipeSectionState}
  \uses{def:physics:phyx-mini-0724:phyxminiproblems-problemphyxmini0724-lengthquantity}
  The physical state at one marked cross section. In particular, the pressure and speed at section 2 are independent observables; neither is defined from the recorded answer
\end{definition}
\begin{proof}
  This is the declared model definition of Pipe Section State.
\end{proof}
\begin{definition}[Rising Pipe Setup]
  \label{def:physics:phyx-mini-0724:phyxminiproblems-problemphyxmini0724-risingpipesetup}
  \lean{PhyXMiniProblems.ProblemPhyXMini0724.RisingPipeSetup}
  \uses{def:physics:phyx-mini-0724:phyxminiproblems-problemphyxmini0724-lengthquantity, def:physics:phyx-mini-0724:phyxminiproblems-problemphyxmini0724-massdensityquantity, def:physics:phyx-mini-0724:phyxminiproblems-problemphyxmini0724-accelerationquantity, def:physics:phyx-mini-0724:phyxminiproblems-problemphyxmini0724-pipesection, def:physics:phyx-mini-0724:phyxminiproblems-problemphyxmini0724-fluidkind, def:physics:phyx-mini-0724:phyxminiproblems-problemphyxmini0724-pressurereference, def:physics:phyx-mini-0724:phyxminiproblems-problemphyxmini0724-waterpipefigure, def:physics:phyx-mini-0724:phyxminiproblems-problemphyxmini0724-pipesectionstate}
  The independent physical quantities and ideal-flow descriptors for the experiment. verticalRise is the centreline rise from section 1 to section 2, as indicated by the vertical 2.0 m dimension arrow
\end{definition}
\begin{proof}
  This is the declared model definition of Rising Pipe Setup.
\end{proof}
\begin{definition}[Matches Ideal Water Pipe Scenario]
  \label{def:physics:phyx-mini-0724:phyxminiproblems-problemphyxmini0724-matchesidealwaterpipescenario}
  \lean{PhyXMiniProblems.ProblemPhyXMini0724.MatchesIdealWaterPipeScenario}
  \uses{def:physics:phyx-mini-0724:phyxminiproblems-problemphyxmini0724-risingpipesetup}
  The textbook Bernoulli model used for the stated water-flow scenario. Gauge pressures share the same atmospheric reference, so their difference may be used in Bernoulli's equation
\end{definition}
\begin{proof}
  This is the declared model definition of Matches Ideal Water Pipe Scenario.
\end{proof}
\begin{definition}[Matches Primary Pipe Figure]
  \label{def:physics:phyx-mini-0724:phyxminiproblems-problemphyxmini0724-matchesprimarypipefigure}
  \lean{PhyXMiniProblems.ProblemPhyXMini0724.MatchesPrimaryPipeFigure}
  \uses{def:physics:phyx-mini-0724:phyxminiproblems-problemphyxmini0724-risingpipesetup}
  Transcription of the literal labels, section ordering, green arrows, gauge taps, diameter markers, and rising connector in the primary raster. This contains no numerical value for the upper gauge pressure
\end{definition}
\begin{proof}
  This is the declared model definition of Matches Primary Pipe Figure.
\end{proof}
\begin{definition}[Matches Problem Readouts]
  \label{def:physics:phyx-mini-0724:phyxminiproblems-problemphyxmini0724-matchesproblemreadouts}
  \lean{PhyXMiniProblems.ProblemPhyXMini0724.MatchesProblemReadouts}
  \uses{def:physics:phyx-mini-0724:phyxminiproblems-problemphyxmini0724-lengthinmeters, def:physics:phyx-mini-0724:phyxminiproblems-problemphyxmini0724-lengthincentimeters, def:physics:phyx-mini-0724:phyxminiproblems-problemphyxmini0724-speedinmeterspersecond, def:physics:phyx-mini-0724:phyxminiproblems-problemphyxmini0724-pressureinkilopascals, def:physics:phyx-mini-0724:phyxminiproblems-problemphyxmini0724-risingpipesetup}
  The five numerical values supplied by the prose and the primary raster. The upper speed and upper pressure deliberately do not occur as numerical data
\end{definition}
\begin{proof}
  This is the declared model definition of Matches Problem Readouts.
\end{proof}
\begin{definition}[Uses Textbook Water Density And Gravity]
  \label{def:physics:phyx-mini-0724:phyxminiproblems-problemphyxmini0724-usestextbookwaterdensityandgravity}
  \lean{PhyXMiniProblems.ProblemPhyXMini0724.UsesTextbookWaterDensityAndGravity}
  \uses{def:physics:phyx-mini-0724:phyxminiproblems-problemphyxmini0724-densityinkilogramspercubicmeter, def:physics:phyx-mini-0724:phyxminiproblems-problemphyxmini0724-accelerationinmeterspersecondsquared, def:physics:phyx-mini-0724:phyxminiproblems-problemphyxmini0724-risingpipesetup}
  Standard textbook calibration implicit in the numerical answer: water has density 1000 kg/m³, and gravitational acceleration is taken as 9.8 m/s²
\end{definition}
\begin{proof}
  This is the declared model definition of Uses Textbook Water Density And Gravity.
\end{proof}
\begin{definition}[Satisfies Circular Rising Pipe Geometry]
  \label{def:physics:phyx-mini-0724:phyxminiproblems-problemphyxmini0724-satisfiescircularrisingpipegeometry}
  \lean{PhyXMiniProblems.ProblemPhyXMini0724.SatisfiesCircularRisingPipeGeometry}
  \uses{def:physics:phyx-mini-0724:phyxminiproblems-problemphyxmini0724-lengthinmeters, def:physics:phyx-mini-0724:phyxminiproblems-problemphyxmini0724-areainsquaremeters, def:physics:phyx-mini-0724:phyxminiproblems-problemphyxmini0724-risingpipesetup}
  Both labelled diameters describe circular cross sections, and the marked 2.0 m arrow is the centreline elevation difference. These are geometric relations, not solved pressure or velocity formulas
\end{definition}
\begin{proof}
  This is the declared model definition of Satisfies Circular Rising Pipe Geometry.
\end{proof}
\begin{definition}[Satisfies Ideal Pipe Flow Laws]
  \label{def:physics:phyx-mini-0724:phyxminiproblems-problemphyxmini0724-satisfiesidealpipeflowlaws}
  \lean{PhyXMiniProblems.ProblemPhyXMini0724.SatisfiesIdealPipeFlowLaws}
  \uses{def:physics:phyx-mini-0724:phyxminiproblems-problemphyxmini0724-lengthinmeters, def:physics:phyx-mini-0724:phyxminiproblems-problemphyxmini0724-areainsquaremeters, def:physics:phyx-mini-0724:phyxminiproblems-problemphyxmini0724-speedinmeterspersecond, def:physics:phyx-mini-0724:phyxminiproblems-problemphyxmini0724-pressureinpascals, def:physics:phyx-mini-0724:phyxminiproblems-problemphyxmini0724-densityinkilogramspercubicmeter, def:physics:phyx-mini-0724:phyxminiproblems-problemphyxmini0724-accelerationinmeterspersecondsquared, def:physics:phyx-mini-0724:phyxminiproblems-problemphyxmini0724-risingpipesetup}
  The steady incompressible continuity equation and Bernoulli equation between the two marked section centres, written in coherent SI readouts. The latter uses a common gauge-pressure reference and contains no problem-specific solved upper pressure
\end{definition}
\begin{proof}
  This is the declared model definition of Satisfies Ideal Pipe Flow Laws.
\end{proof}
\begin{lemma}[upper speed from continuity]
  \label{lem:physics:phyx-mini-0724:phyxminiproblems-problemphyxmini0724-upper-speed-from-continuity}
  \lean{PhyXMiniProblems.ProblemPhyXMini0724.upper_speed_from_continuity}
  \uses{def:physics:phyx-mini-0724:phyxminiproblems-problemphyxmini0724-speedinmeterspersecond, def:physics:phyx-mini-0724:phyxminiproblems-problemphyxmini0724-risingpipesetup, def:physics:phyx-mini-0724:phyxminiproblems-problemphyxmini0724-matchesproblemreadouts, def:physics:phyx-mini-0724:phyxminiproblems-problemphyxmini0724-satisfiescircularrisingpipegeometry, def:physics:phyx-mini-0724:phyxminiproblems-problemphyxmini0724-satisfiesidealpipeflowlaws}
  Circular-section geometry and continuity give v₂ = 5 (6/4)² = 45/4 m/s
\end{lemma}
\begin{proof}
  The conclusion follows from the governing relations and figure readouts named in the statement of upper speed from continuity.
\end{proof}
\begin{lemma}[upper gauge pressure in pascals]
  \label{lem:physics:phyx-mini-0724:phyxminiproblems-problemphyxmini0724-upper-gauge-pressure-in-pascals}
  \lean{PhyXMiniProblems.ProblemPhyXMini0724.upper_gauge_pressure_in_pascals}
  \uses{def:physics:phyx-mini-0724:phyxminiproblems-problemphyxmini0724-pressureinpascals, def:physics:phyx-mini-0724:phyxminiproblems-problemphyxmini0724-risingpipesetup, def:physics:phyx-mini-0724:phyxminiproblems-problemphyxmini0724-matchesproblemreadouts, def:physics:phyx-mini-0724:phyxminiproblems-problemphyxmini0724-usestextbookwaterdensityandgravity, def:physics:phyx-mini-0724:phyxminiproblems-problemphyxmini0724-satisfiescircularrisingpipegeometry, def:physics:phyx-mini-0724:phyxminiproblems-problemphyxmini0724-satisfiesidealpipeflowlaws}
  Substitution into Bernoulli's equation gives the exact idealized upper gauge pressure 18475/4 Pa = 739/160 kPa
\end{lemma}
\begin{proof}
  The conclusion follows from the governing relations and figure readouts named in the statement of upper gauge pressure in pascals.
\end{proof}
\begin{definition}[Answer Choice]
  \label{def:physics:phyx-mini-0724:phyxminiproblems-problemphyxmini0724-answerchoice}
  \lean{PhyXMiniProblems.ProblemPhyXMini0724.AnswerChoice}
  Labels of the four pressure readings printed with the problem
\end{definition}
\begin{proof}
  This is the declared model definition of Answer Choice.
\end{proof}
\begin{definition}[displayed Gauge Pressure Kilopascals]
  \label{def:physics:phyx-mini-0724:phyxminiproblems-problemphyxmini0724-displayedgaugepressurekilopascals}
  \lean{PhyXMiniProblems.ProblemPhyXMini0724.displayedGaugePressureKilopascals}
  \uses{def:physics:phyx-mini-0724:phyxminiproblems-problemphyxmini0724-answerchoice}
  Kilopascal value printed beside each displayed answer choice
\end{definition}
\begin{proof}
  This is the declared model definition of displayed Gauge Pressure Kilopascals.
\end{proof}
\begin{definition}[Rounds To Nearest Tenth Kilopascal]
  \label{def:physics:phyx-mini-0724:phyxminiproblems-problemphyxmini0724-roundstonearesttenthkilopascal}
  \lean{PhyXMiniProblems.ProblemPhyXMini0724.RoundsToNearestTenthKilopascal}
  A reading displayed to the nearest 0.1 kPa differs from the exact idealized readout by strictly less than 0.05 kPa
\end{definition}
\begin{proof}
  This is the declared model definition of Rounds To Nearest Tenth Kilopascal.
\end{proof}
\begin{definition}[Matches Displayed Upper Gauge Pressure]
  \label{def:physics:phyx-mini-0724:phyxminiproblems-problemphyxmini0724-matchesdisplayeduppergaugepressure}
  \lean{PhyXMiniProblems.ProblemPhyXMini0724.MatchesDisplayedUpperGaugePressure}
  \uses{def:physics:phyx-mini-0724:phyxminiproblems-problemphyxmini0724-pressureinkilopascals, def:physics:phyx-mini-0724:phyxminiproblems-problemphyxmini0724-risingpipesetup, def:physics:phyx-mini-0724:phyxminiproblems-problemphyxmini0724-answerchoice, def:physics:phyx-mini-0724:phyxminiproblems-problemphyxmini0724-displayedgaugepressurekilopascals, def:physics:phyx-mini-0724:phyxminiproblems-problemphyxmini0724-roundstonearesttenthkilopascal}
  A displayed choice agrees with the modeled upper gauge to its precision
\end{definition}
\begin{proof}
  This is the declared model definition of Matches Displayed Upper Gauge Pressure.
\end{proof}
\begin{definition}[Is Unique Matching Upper Gauge Pressure]
  \label{def:physics:phyx-mini-0724:phyxminiproblems-problemphyxmini0724-isuniquematchinguppergaugepressure}
  \lean{PhyXMiniProblems.ProblemPhyXMini0724.IsUniqueMatchingUpperGaugePressure}
  \uses{def:physics:phyx-mini-0724:phyxminiproblems-problemphyxmini0724-risingpipesetup, def:physics:phyx-mini-0724:phyxminiproblems-problemphyxmini0724-answerchoice, def:physics:phyx-mini-0724:phyxminiproblems-problemphyxmini0724-matchesdisplayeduppergaugepressure}
  Exactly one printed choice agrees with the modeled upper gauge reading
\end{definition}
\begin{proof}
  This is the declared model definition of Is Unique Matching Upper Gauge Pressure.
\end{proof}
% --- Archon declaration topology end ---
% --- Archon physics formalization source end ---
